\section{Scalar Products}
bx{
\en{
    \item 
    Let $A$ be a symmetric $n \times n$ matrix. 
    Given two column vectors $X, Y \in \mathbf{R}^n$, define 
    \begin{equation*}
        \langle X, Y \rangle = {^{t\!}X} A Y. 
    \end{equation*}
    \ea{
        \item 
        Show that this symbol satisfies the first three properties of a scalar product. 
        \item 
        Given an example of a $1 \times 1$ matrix and a non-zero $2 \times 2$ matrix such that the fourth property is not satisfied. 
        If this fourth property is satisfied, that is $^{t\!}XAX > 0$ for all $X \neq O$, then the matrix $A$ is called \textbf{positive definite}. 
        \item 
        Given an example of $2 \times 2$ matrix which is symmetric and positive definite. 
        \item 
        Let $a > 0$, and let 
        \begin{equation*}
            A = 
            \begin{pmatrix}
                a & b \\
                b & d
            \end{pmatrix}. 
        \end{equation*}
        Prove that $A$ is positive definite if and only if $ad - b^2 > 0$. 
        [\textit{Hint}: Let $X = {^{t\!}(x,y)}$ and complete the square in the expression $^{t\!}X A X$.] 
        \item 
        If $a < 0$ show that $A$ is not positive definite. 
    }
    \begin{Solution}
        \ea{
            \item 
            \begin{proof}
                We first observe the result of the scalar product we defined is a $1 \times 1$ matrix, which means that it is simply a scalar. 

                We begin by showing the first property. 
                We know 
                \begin{equation*}
                    \langle X, Y \rangle = {^{t\!}X} A Y \quad \text{and} \quad \langle Y, X \rangle = {^{t\!}Y} A X. 
                \end{equation*}
                And we observe that 
                \begin{equation*}
                    {^{t\!}}({^{t\!}X} A Y) = {^{t\!}Y} {^{t\!\!}A} X = {^{t\!}Y} A X
                \end{equation*}
                because $A$ is symmetric. 
                Then we have 
                \begin{equation*}
                    \langle X, Y \rangle = {^{t\!}}\langle X, Y \rangle = \langle Y, X \rangle 
                \end{equation*}
                because $\langle X, Y \rangle$ and $\langle Y, X \rangle$ are both scalars, which proving the first property. 

                We continue to show the second property. 
                Let $Z \in \mathbf{R}^n$ is a column vector. 
                Then by the distributivity of matrices, we have 
                \begin{equation*}
                    \langle X, Y + Z \rangle = {^{t\!}X} A (Y + Z) = {^{t\!}X} A Y + {^{t\!}X} A Z = \langle X, Y \rangle + \langle X, Z \rangle, 
                \end{equation*}
                which proves the second property. 

                We finally come to the third property. 
                If $x$ is a number, then 
                \begin{equation*}
                    \langle xX, Y \rangle = {^{t\!}(xX)} A Y = x {^{t\!}(xX)} A Y = x \langle X, Y\rangle. 
                \end{equation*}
                On the other hand, we have 
                \begin{equation*}
                    \langle X, xY \rangle = {^{t\!}X} A (xY) = x {^{t\!}X} A Y = x \langle X, Y \rangle. 
                \end{equation*}
                Thus we conclude the whole proof. 
            \end{proof}
            \item 
            $\begin{pmatrix} -1 \end{pmatrix}$, $\begin{pmatrix} -1 & 0 \\ 0 & -1 \end{pmatrix}$ 
            \item
            $\begin{pmatrix} 1 & 0 \\ 0 & 1 \end{pmatrix}$
            \item 
            \begin{proof}
                Let $X = {^{t\!}(x, y)}$. 
                Then 
                \begin{equation*}
                    \langle X, X \rangle = ax^2 + 2bxy + dy^2 = a \left( x + \frac{b}{a}y \right)^2 + \frac{(ad - b^2)}{a} y^2. 
                \end{equation*}
                If $ad - b^2 > 0$, then $\langle X, X \rangle$ is a sum of squares with positive coefficients, and one of the two terms is not $0$, so $\langle X, X \rangle > 0$. 
                If $ad - b^2 \leq 0$, then give $y$ any non-zero value, and let $x = -by/a$. 
                Then $\langle X, X \rangle \leq 0$. 
            \end{proof}
            \item 
            \begin{proof}
                Let $X = {^{t\!}(1, 0)} \neq O$ (of course). 
                Then 
                \begin{equation*}
                    \langle X, X \rangle = a < 0, 
                \end{equation*}
                which contradicts the definition of the positive definite matrix. 
                Hence $a < 0$ implies $A$ is not positive definite. 
        \end{proof}
        }
    \end{Solution}
    \item
    \item
    \item
    Let $A$ be an $n \times n$ matrix. 
    Define the \textbf{trace} of $A$ to be the sum of the diagonal elements. 
    Thus if $A = (a_{ij})$, then 
    \begin{equation*}
        \tr(A) = \sum_{i = 1}^n a_{ii}. 
    \end{equation*}
    For instance, if 
    \begin{equation*}
        A = 
        \begin{pmatrix}
            1 & 2 \\
            3 & 4
        \end{pmatrix}, 
    \end{equation*}
    then $\tr(A) = 1 + 4 = 5$. 
    If 
    \begin{equation*}
        A = 
        \begin{pmatrix}
            1 &-1 & 5 \\
            2 & 1 & 3 \\
            1 &-4 & 7
        \end{pmatrix}, 
    \end{equation*}
    then $\tr(A) = 9$. 
    \item
    \item
    If $A$ is a symmetric square matrix, show that $\tr(AA) \geq 0$, and $= 0$ if and only if $A = O$. 
    \begin{Solution}
        \begin{proof}
            Let $AA = (c_{ij})$. 
            Then 
            \begin{equation*}
                c_{ii} = \sum_{k = 1}^{n} a_{ik} a_{ki}. 
            \end{equation*}
            Hence 
            \begin{equation*}
                \tr(AA) = \sum_{i = 1}^{n} \sum_{k = 1}^{n} a_{ik} a_{ki}. 
            \end{equation*}
            and $a_{ik} a_{ki} = a_{ik}^2$ because $A$ is symmetric, so the trace is a sum of squares, hence $\geq 0$. 
        \end{proof}
    \end{Solution}
    \item
    \item
    \ea{
        \item 
        Prove in general that if $A, B$ are square $n \times n$ matrices, then 
        \begin{equation*}
            \tr(AB) = \tr(BA). 
        \end{equation*}
        \item 
        If $C$ is an $n \times n$ matrix which has an inverse, then $\tr(C^{-1}AC) = \tr(A)$. 
    }
    \begin{Solution}
        \ea{
            \item 
            \begin{proof}
                \begin{equation*}
                    \tr(AB) = \sum_i \sum_j a_{ij} b_{ji} = \sum_j \sum_i b_{ji} a_{ij}
                \end{equation*}
                But any pair of letters can be used to denote indices in this sum, which can be rewritten more neutrally. 
                \begin{equation*}
                    \sum_{r = 1}^{n} \sum_{s=1}^n b_{rs}a_{sr}. 
                \end{equation*}
                This is precisely the trace $\tr(BA)$. 
            \end{proof}
            \item 
            $\tr(C^{-1}AC) = \tr(ACC^{-1})$ by part (a), so $= \tr(A)$. 
        }
    \end{Solution}
}}
